\heading{热学-18春-期中回忆}

\noteinfo{题目来源于赛艇，答案由AI生成。第 14、16 题缺失。}

\textbf{一、选择题（回忆到的部分）}
\begin{question}
节流过程中熵的变化。
\begin{choiceoptions}
\item 增大
\item 减小
\item 不变
\item 不确定
\end{choiceoptions}
\begin{solution}
节流是不可逆过程，理想情形下焓守恒，但熵增不减。
\ansmath{\text{A}}
\end{solution}
\end{question}

\begin{question}
（回忆题）在固定压强下给物质以热量，哪种物质可能不会改变温度？
\begin{choiceoptions}
\item 理想气体
\item 满足麦克斯韦速度分布的气体
\item 处于相变点的气体
\item 单原子气体
\end{choiceoptions}
\begin{solution}
与 17 春补录题相同，处于相变点时热量可全部用于相变。
\ansmath{\text{C}}
\end{solution}
\end{question}

\begin{question}
（回忆题）测温属性与温度之间应满足何种关系？
\begin{choiceoptions}
\item 线性函数
\item 升函数
\item 降函数
\item 单调函数
\end{choiceoptions}
\begin{solution}
只要求温度与测温属性一一对应，故应为单调函数。
\ansmath{\text{D}}
\end{solution}
\end{question}

\begin{question}
关于麦克斯韦速率分布，下列哪种说法正确？
\begin{choiceoptions}
\item 是高斯分布的一种形式
\item 是热平衡条件下的一类特定速率分布
\item 在存在外势场时一定失效
\item 在任何条件下都成立
\end{choiceoptions}
\begin{solution}
速率分布本身不是高斯分布，也不是任意条件下都成立；在热平衡下它描述平衡气体的一类确定分布律。
\ansmath{\text{B}}
\end{solution}
\end{question}

\begin{question}
纯物质气液临界点的表面张力系数。
\begin{choiceoptions}
\item 大于 0
\item 小于 0
\item 等于 0
\item 不确定
\end{choiceoptions}
\begin{solution}
临界点附近液、气两相趋于不可区分，因此界面张力消失。
\ansmath{\text{C}}
\end{solution}
\end{question}

\textbf{二、简答题}
\begin{question}
请写出五个热力学态函数。
\begin{solution}
常见的五个态函数是：内能 $U$、焓 $H$、熵 $S$、Helmholtz 自由能 $F$、Gibbs 自由能 $G$。
\anstext{$U,\ H,\ S,\ F,\ G$}
\end{solution}
\end{question}

\begin{question}
冰箱的工作物质能不能用理想气体？
\begin{solution}
冰箱需要利用工质的液化与汽化过程来进行大规模吸放热，而理想气体模型没有真实相变，因此不能作为实际制冷工质模型。
\anstext{一般不能。因为制冷循环依赖真实工质的相变潜热。}
\end{solution}
\end{question}

\begin{question}
书 P277 图 6.15 所示：起始时是一个大肥皂泡和一个小肥皂泡，描述之后的现象。
\begin{solution}
由于小泡内压更大，气体会从小泡流向大泡，于是小泡继续变小甚至消失，大泡继续变大。
\anstext{小泡变小，大泡变大。}
\end{solution}
\end{question}

\begin{question}
一级相变的两个特征。
\begin{solution}
一级相变沿相变线具有潜热，同时比容（或密度）出现跃变。
\anstext{具有潜热；体积（或密度）有跃变。}
\end{solution}
\end{question}

\begin{question}
描述什么是过热液体和过冷蒸气，并在图中表示分别是等温线的哪一段。
\begin{solution}
过热液体是温度高于该压强下平衡沸点、但仍保持液态的亚稳液体；过冷蒸气是温度低于该压强下饱和温度、但仍保持气态的亚稳蒸气。它们都对应等温线进入两相共存区附近的亚稳延拓段。
\anstext{分别对应液相支和气相支向两相区延拓的亚稳段。}
\end{solution}
\end{question}

\textbf{三、解答题（按回忆稿整理）}
\begin{question}
算奥托循环的效率。
\begin{solution}
对理想气体 Otto 循环，效率仅与压缩比 $r$ 和绝热指数 $\gamma$ 有关：
\[
\eta=1-\frac{1}{r^{\gamma-1}}.
\]
\ansmath{\eta_{\mathrm{Otto}}=1-r^{1-\gamma}}
\end{solution}
\end{question}

\begin{question}
燃料电池：已知各个键能及水的标准生成焓，求生成一摩尔水释放的热量和消耗的电能。
\begin{solution}
这类题的标准处理是先由键能或生成焓求反应焓变 $\Delta H$，再用电池可逆输出功对应的 Gibbs 自由能变 $\Delta G$。若只问“释放热量”，常取定压反应放热
\[
Q_p=-\Delta H.
\]
若问消耗的电能（或最多输出的电功），则取
\[
W_{\mathrm{elec,max}}=-\Delta G.
\]
\anstext{放热量由 $-\Delta H$ 给出；可逆电功由 $-\Delta G$ 给出。}
\end{solution}
\end{question}

\begin{question}
书 P211：混合前后熵变。
\begin{solution}
若为两份理想气体等温混合，则总熵增写成
\[
\Delta S=-\sum_i n_iR\ln x_i>0.
\]
若是同种气体混合，则无额外混合熵。题目在回忆稿中仅保留了课本页码，故此处记下通用结论。
\ansmath{\Delta S_{\mathrm{mix}}=-\sum_i n_iR\ln x_i\ (\text{异种理想气体})}
\end{solution}
\end{question}

\begin{question}
书 P297 习题 6.32。
\begin{solution}
回忆稿仅给出了课本编号，未保留完整题面，故这里只保留引用信息，待补充原题后可继续誊写。
\anstext{题面缺失，待补。}
\end{solution}
\end{question}

\begin{question}
用范德瓦尔斯气体证明
\[
\left(\frac{\partial U}{\partial V}\right)_T=T\left(\frac{\partial p}{\partial T}\right)_V-p.
\]
\begin{solution}
从热力学基本关系
\[
dU=T\,dS-p\,dV
\]
和 Maxwell 关系
\[
\left(\frac{\partial S}{\partial V}\right)_T=\left(\frac{\partial p}{\partial T}\right)_V
\]
直接得到
\[
\left(\frac{\partial U}{\partial V}\right)_T=T\left(\frac{\partial p}{\partial T}\right)_V-p.
\]
对范德瓦尔斯气体
\[
p=\frac{RT}{V_m-b}-\frac{a}{V_m^2}
\]
代入可进一步化成
\[
\left(\frac{\partial U}{\partial V}\right)_T=\frac{a}{V_m^2}.
\]
\ansmath{\left(\frac{\partial U}{\partial V}\right)_T=T\left(\frac{\partial p}{\partial T}\right)_V-p}
\end{solution}
\end{question}

\begin{question}
书 P352 习题 7.10。
\begin{solution}
回忆稿仅保留了课本编号，缺少完整题面，因此此处保留空位。
\anstext{题面缺失，待补。}
\end{solution}
\end{question}

\begin{question}
一永动机如图：上端左端都开口，除 $h_1$ 外其余几何量、液体密度、管径 $d$、表面张力系数 $\sigma$ 已知。
\begin{enumerate}[label=（\arabic*）,leftmargin=2.8em,itemsep=0.2em]
\item 求 $h_1$；
\item 求侧面那个管口液面形状（弯曲液面角度）；
\item 说明无论怎么取值，该永动机都不可能工作。
\end{enumerate}
\begin{solution}
该题的核心是毛细压强与静压平衡：若两侧自由液面都与外界相通，则任一闭合回路上的压强增量代数和必须为零。毛细弯月面只能重新分配局部压差，不能在一周内产生净压头。

因此
\begin{enumerate}[label=（\arabic*）,leftmargin=2.8em,itemsep=0.2em]
\item $h_1$ 由 Laplace 压差与静水压差平衡确定，形式上满足 $\rho g h_1=2\sigma/R$（圆管中 $R$ 由接触角决定）。
\item 侧管液面形状由 Young--Laplace 方程和接触角条件共同确定。
\item 无论怎样选取参数，沿闭合回路总机械能增益都为零，因此该装置不可能持续输出功。
\end{enumerate}
\anstext{永动机不可能工作；毛细现象不能提供闭路净压头。}
\end{solution}
\end{question}


